\chapter{Orchestration with Snowflake Tasks and Task Graphs}
\index{tasks}\index{task graph}\index{orchestration}\index{serverless tasks}\index{Week 8}%
\begin{objectives}
\begin{itemize}
    \item Schedule tasks with CRON expressions or interval schedules and configure execution parameters.
    \item Compose multi-step task trees (DAGs) with root tasks and \texttt{AFTER} dependencies.
    \item Choose between warehouse-backed and serverless tasks and explain the billing difference.
    \item Wire error notifications and failure suspension so task graphs fail loudly and stop cascading.
    \item Write idempotent task steps so retries never double-apply work.
    \item Architect a serverless, self-healing ELT orchestration engine for 50 dependent business transformations.
\end{itemize}
\end{objectives}

\begin{crosscloud}
Every cloud grew a scheduler, a DAG engine, and a serverless variant; Snowflake tasks are simply the in-database member of that family.

\vspace{2mm}
{\footnotesize
\begin{tabular}{@{}p{2.8cm}p{3.0cm}p{3.0cm}p{3.0cm}@{}}
\toprule
\textbf{Capability} & \textbf{AWS} & \textbf{Azure} & \textbf{GCP} \\
\midrule
Cron scheduler & EventBridge Scheduler & Azure Functions timer / Logic Apps & Cloud Scheduler \\
DAG orchestrator & Step Functions / MWAA (Airflow) & Data Factory pipelines & Cloud Composer (Airflow) \\
In-warehouse scheduling & Redshift scheduled queries & Synapse/ADF triggers & BigQuery scheduled queries \\
Serverless per-run billing & Lambda-style per invocation & Consumption plans & Per-request / slot-time \\
Failure notification & SNS & Azure Monitor / Logic Apps & Cloud Monitoring + Pub/Sub \\
\addlinespace
Snowflake equivalent & \multicolumn{3}{l}{Tasks (CRON or interval), task graphs via \texttt{AFTER}, serverless compute, \texttt{ERROR\_INTEGRATION}} \\
\bottomrule
\end{tabular}}

\vspace{1mm}
\textbf{Microsoft Fabric} orchestrates with Data Factory pipelines and notebook scheduling; \textbf{MongoDB} relies on Atlas triggers and external schedulers. Snowflake's edge is proximity: when the entire transformation lives in SQL, orchestrating inside the database removes a whole external system---and its credentials, networking, and failure modes.
\end{crosscloud}

\section{Week 8 Roadmap}
Chapters 6 and 7 produced continuous inputs and change deltas. Something must still \emph{decide when work runs}, \emph{in what order}, and \emph{what happens when it fails}. Snowflake tasks answer inside the platform, replacing a surprising share of external orchestrators for ELT workloads \cite{snowflakeTasksDocs}.

Monday covers scheduling and execution parameters. Tuesday builds task trees with \texttt{AFTER}. Wednesday addresses serverless tasks, error notifications, and failure containment. Thursday constructs a gated three-task CDC graph. Friday scales to a 50-transformation self-healing orchestration engine.

\begin{keyterms}
\textbf{Task}: a scheduled Snowflake object executing SQL or a stored procedure call.\quad
\textbf{Root task}: the scheduled head of a task graph; child tasks hang below it.\quad
\textbf{Task tree / DAG}: a dependency graph of tasks linked by \texttt{AFTER} clauses.\quad
\textbf{Serverless task}: a task run on Snowflake-managed compute, billed per second of usage.\quad
\textbf{ERROR\_INTEGRATION}: a notification integration that emits task failure events to SNS, queues, or email.
\end{keyterms}

\section{Monday: Scheduling and Execution Parameters}
\index{task!scheduling}\index{CRON}
A task binds four things: a schedule, a compute choice, optional preconditions, and a body (a single SQL statement, a multi-statement script, or a stored procedure call). Schedules come in two forms: \texttt{SCHEDULE = 'n MINUTE'} intervals, and full \texttt{USING CRON} expressions with an explicit time zone for calendar-shaped timing such as ``06:00 Europe/Berlin on weekdays.''

\begin{lstlisting}[language=SQL, caption={Two scheduling styles and execution parameters}]
CREATE OR REPLACE TASK t_daily_billing
  WAREHOUSE = etl_wh
  SCHEDULE = 'USING CRON 0 6 * * MON-FRI Europe/Berlin'
  AS CALL billing.sp_build_daily_charges();

CREATE OR REPLACE TASK t_incremental
  WAREHOUSE = etl_wh
  SCHEDULE = '15 MINUTE'
  USER_TASK_TIMEOUT_MS = 600000          -- fail fast, not silently hung
  SUSPEND_TASK_AFTER_NUM_FAILURES = 3     -- stop cascading failure
  AS CALL elt.sp_merge_orders();

-- New tasks are created suspended
ALTER TASK t_daily_billing RESUME;
\end{lstlisting}

Two governance facts deserve emphasis. First, tasks are created \emph{suspended}; forgetting \texttt{RESUME} is the most common first-day bug. Second, \texttt{USER\_TASK\_TIMEOUT\_MS} and \texttt{SUSPEND\_TASK\_AFTER\_NUM\_FAILURES} are not optional hardening---they are the difference between a stuck task burning credits for a weekend and a contained incident.

\begin{pitfall}
\textbf{Resuming a child task before its root} silently leaves the tree inert: child tasks resume into a waiting state and run only when the root fires and the dependency chain above them has succeeded. Operational runbooks should resume trees leaf-last, root-first---or script the resume order.
\end{pitfall}

\section{Tuesday: Task Trees and DAGs}
\index{task tree}\index{AFTER clause}\index{DAG}
A task graph forms when child tasks declare \texttt{AFTER <parent>}. Exactly one root carries the schedule; children inherit the graph's rhythm and run only after every declared parent succeeds in a given run. Dependencies can fan out (one parent, many children) and fan in (a child with multiple \texttt{AFTER} parents), producing arbitrary DAGs.

\begin{lstlisting}[language=SQL, caption={A gated CDC task tree}]
CREATE OR REPLACE TASK t_gate
  WAREHOUSE = etl_wh
  SCHEDULE = '5 MINUTE'
  AS CALL elt.sp_noop();   -- anchor; children carry the guard

CREATE OR REPLACE TASK t_merge_customers
  AFTER t_gate
  WHEN SYSTEM$STREAM_HAS_DATA('raw.customer_stream')
  AS CALL elt.sp_merge_customers();

CREATE OR REPLACE TASK t_merge_orders
  AFTER t_gate
  WHEN SYSTEM$STREAM_HAS_DATA('raw.order_stream')
  AS CALL elt.sp_merge_orders();

CREATE OR REPLACE TASK t_cleanup
  AFTER t_merge_customers, t_merge_orders
  AS CALL elt.sp_truncate_staging();
\end{lstlisting}

The \texttt{WHEN} clause deserves careful reading: it is a per-run precondition evaluated when the task's parents complete. If the predicate is false, the task \emph{skips}---and, critically, its children skip too. That makes a gated parent the standard ``don't run the graph when there's no data'' pattern, saving credits across an entire subtree.

\begin{tipbox}
Keep the root task a trivial anchor and push all gating logic into \texttt{WHEN} clauses on first-level children. Graphs whose root does heavy work are hard to re-point at new schedules without disturbing the tree.
\end{tipbox}

\subsection{Execution Semantics and Overlap}
\index{task!overlap}
A graph run that exceeds its schedule interval does not overlap itself: the next run starts only after the current one completes. Latency-sensitive graphs must therefore be designed so worst-case runtime fits inside the interval, or split so slow branches live on their own roots. \texttt{ALTER TASK ... SET TARGET\_COMPLETION\_INTERVAL} and query history make this measurable rather than assumed.

\section{Wednesday: Serverless Tasks and Error Handling}
\index{serverless task}\index{error notification}\index{ERROR\_INTEGRATION}
Omitting \texttt{WAREHOUSE} makes a task \textbf{serverless}: Snowflake manages and autosizes the compute, billing per second of task compute. Serverless is the right default for short, frequent, spiky steps---a five-minute MERGE every five minutes leaves a warehouse-backed design paying for idle time between runs---while long, steady, or resource-hungry steps still justify a dedicated warehouse whose size you control.

\subsection{Failure Containment}
\index{SUSPEND\_TASK\_AFTER\_NUM\_FAILURES}
Three mechanisms convert task failures from silent data rot into actionable incidents:
\begin{enumerate}
    \item \textbf{\texttt{ERROR\_INTEGRATION}} wires task failures to a notification integration (SNS to Slack or PagerDuty, or email), so the on-call engineer learns about the failure in minutes.
    \item \textbf{\texttt{SUSPEND\_TASK\_AFTER\_NUM\_FAILURES}} halts a repeatedly failing task before it cascades: a root suspended after three failures stops the whole subtree rather than corrupting downstream layers.
    \item \textbf{Idempotent steps} make retry safe. Every task body should be re-runnable: \texttt{MERGE} with keys, \texttt{INSERT OVERWRITE}, or swap-based rebuilds. A task whose retry double-applies work converts a transient failure into a data-quality incident.
\end{enumerate}

\begin{notebox}
Task observability lives in \texttt{INFORMATION\_SCHEMA.TASK\_HISTORY} and \texttt{SNOWFLAKE.ACCOUNT\_USAGE.TASK\_HISTORY}: run states, error messages, durations, and skip reasons. Build the operations dashboard on day one, not after the first missed SLA.
\end{notebox}

\section{Thursday Hands-On Project: A Gated Three-Task CDC Graph}
\index{hands-on project!tasks}
This project assembles Chapters 6--7 machinery into a self-managing graph: gate on stream data, merge deltas, clean staging.

\subsection*{Scenario}
The retailer from Chapter 7 wants its MySQL-derived streams consumed every five minutes, but only when data exists, with staging truncated after successful application and on-call paged on failure.

\subsection*{Procedure}
\begin{enumerate}
    \item Reuse \texttt{raw.customer\_stream} and the \texttt{MERGE} logic from Chapter 7, wrapping the merge in \texttt{elt.sp\_merge\_customers()}.
    \item Create the anchor root \texttt{t\_gate} on a 5-minute schedule.
    \item Create \texttt{t\_merge\_customers} with \texttt{AFTER t\_gate} and the \texttt{WHEN SYSTEM\$STREAM\_HAS\_DATA} guard.
    \item Create \texttt{t\_cleanup} after the merge, truncating staging.
    \item Attach an \texttt{ERROR\_INTEGRATION} notification integration and set \texttt{SUSPEND\_TASK\_AFTER\_NUM\_FAILURES = 3} on each task.
    \item Resume in order (children, then root); insert test rows and watch one full run in \texttt{TASK\_HISTORY}.
    \item Force a failure (bad privilege on the procedure); confirm the error notification fires and the task auto-suspends after three failures.
    \item Fix the fault, resume, and prove the retry applied deltas exactly once.
\end{enumerate}

\subsection*{Deliverables}
\begin{itemize}
    \item \textbf{DDL script:} procedures, three tasks, notification integration, resume sequence.
    \item \textbf{Run evidence:} \texttt{TASK\_HISTORY} excerpts showing a skip (no data), a success, and a contained failure.
    \item \textbf{Idempotency proof:} row counts unchanged across the recovery retry.
    \item \textbf{Graph diagram:} dependencies, schedules, and gating predicates.
\end{itemize}

\subsection*{Acceptance Criteria}
\begin{itemize}
    \item No graph run occurs when the stream is empty; \texttt{TASK\_HISTORY} shows the skip reason.
    \item Failures notify through the integration and the faulting task suspends after the configured threshold.
    \item Cleanup runs only after the merge succeeds, demonstrated by an induced merge failure.
    \item Post-fix retry neither loses nor duplicates deltas.
\end{itemize}

\subsection*{Stretch Goals}
\begin{itemize}
    \item Convert the merge tasks to serverless and compare weekly credits against the warehouse-backed version.
    \item Add a second stream-gated branch (orders) and show fan-out/fan-in at cleanup.
    \item Build a \texttt{TASK\_HISTORY} health view: per-task success rate, mean duration, and skip ratio over 7 days.
\end{itemize}

\section{Friday Real-World Architecture: A 50-Transformation Orchestration Engine}
\index{Real-World Architecture!ELT orchestration}
An insurer runs roughly 50 dependent business transformations nightly: reference data, policy facts, claims marts, regulatory extracts. The mandate is serverless, self-healing, and free of external orchestrators.

\subsection{Layered Graph Design}
Fifty tasks in one tree is unmaintainable. The design layers graphs by medallion stage---a bronze root with its children, a silver root gated on bronze completion via a handoff check, and gold marts beneath---with each layer's root scheduled and cross-layer coupling expressed through a small \texttt{layer\_status} control table rather than deep cross-tree \texttt{AFTER} edges. Within a layer, independent marts fan out to maximize parallelism on serverless compute.

\begin{figure}[H]
    \centering
    \resizebox{0.95\textwidth}{!}{%
    \begin{tikzpicture}[
        root/.style={rectangle, rounded corners, draw=primary, thick, fill=primary!8, align=center, minimum width=2.5cm, minimum height=0.9cm, font=\small\bfseries},
        leaf/.style={rectangle, rounded corners, draw=codegreen!60!black, thick, fill=green!8, align=center, minimum width=2.2cm, minimum height=0.8cm, font=\scriptsize},
        ctrl/.style={rectangle, rounded corners, draw=red!60!black, thick, fill=red!6, align=center, minimum width=3.2cm, minimum height=0.9cm, font=\scriptsize\bfseries},
        arrow/.style={-{Stealth[length=2.2mm]}, draw=primary, thick},
        lbl/.style={font=\scriptsize\itshape, text=codegray}
    ]
        \node[root] (broot) at (0, 0) {bronze root\\(CRON 01:00)};
        \node[leaf] (b1) at (-2.6, -1.5) {policies};
        \node[leaf] (b2) at (0, -1.5) {claims};
        \node[leaf] (b3) at (2.6, -1.5) {reference};
        \node[ctrl] (status) at (6.4, -1.5) {layer\_status\\control table};
        \node[root] (sroot) at (6.4, -3.2) {silver root\\(gated on bronze OK)};
        \node[leaf] (s1) at (3.8, -4.7) {policy mart};
        \node[leaf] (s2) at (6.4, -4.7) {claims mart};
        \node[leaf] (s3) at (9.0, -4.7) {reg extracts};
        \draw[arrow] (broot) -- (b1);
        \draw[arrow] (broot) -- (b2);
        \draw[arrow] (broot) -- (b3);
        \draw[arrow] (b1) |- (status);
        \draw[arrow] (b2) -- (status);
        \draw[arrow] (b3) |- (status);
        \draw[arrow] (status) -- node[lbl, right]{all-green gate} (sroot);
        \draw[arrow] (sroot) -- (s1);
        \draw[arrow] (sroot) -- (s2);
        \draw[arrow] (sroot) -- (s3);
    \end{tikzpicture}%
    }
    \caption{Layered orchestration: intra-layer \texttt{AFTER} trees, inter-layer handoff through a control table.}
    \label{fig:chapter8-orchestration}
\end{figure}

\subsection{Self-Healing Properties}
\begin{itemize}
    \item \textbf{Idempotent everything:} each transformation is a keyed \texttt{MERGE} or \texttt{INSERT OVERWRITE}, so any failed branch can simply resume and retry.
    \item \textbf{Bounded blast radius:} \texttt{SUSPEND\_TASK\_AFTER\_NUM\_FAILURES = 3} per task; one poisoned feed stalls its branch, not the night.
    \item \textbf{One notification surface:} a single error integration fans failures into the on-call channel with task name and error text from \texttt{TASK\_HISTORY}.
    \item \textbf{Data-quality gates as tasks:} row-count reconciliation and freshness checks are themselves tasks whose failure suspends publication tasks downstream.
\end{itemize}

\begin{pitfall}
\textbf{An orchestrator without backfill semantics rots.} When a night fails, operators need a sanctioned ``replay layer L for date D'' path---parameterized procedures reading a run-date control table---not fifty manual reruns in Snowsight.
\end{pitfall}

For context on when to outgrow in-platform orchestration, compare Airflow's model---external, Python-defined DAGs with rich cross-system operators---which Chapter 11's dynamic-table discussion revisits from the declarative side \cite{airflowDocs}.

\section{Interview Q\&A}
\begin{enumerate}
    \item \textbf{Q: Which two schedule types can a task use?}\\
    A: A fixed interval (\texttt{'n MINUTE'}) or a \texttt{USING CRON} expression with an explicit time zone.
    \item \textbf{Q: How do child tasks declare dependencies?}\\
    A: With \texttt{AFTER <parent>} clauses; a child runs only when all declared parents succeed in that graph run.
    \item \textbf{Q: What is a serverless task and how is it billed?}\\
    A: A task without a named warehouse, run on Snowflake-managed autosized compute and billed per second of task compute.
    \item \textbf{Q: Why must every task step be idempotent?}\\
    A: Because retries after failure must not double-apply work; keyed MERGEs and INSERT OVERWRITE make re-runs no-ops.
    \item \textbf{Q: How do you alert on a task-graph failure?}\\
    A: Attach \texttt{ERROR\_INTEGRATION} to push failures to SNS/email, and set \texttt{SUSPEND\_TASK\_AFTER\_NUM\_FAILURES} to contain cascades.
\end{enumerate}

\section{Further Reading}
Snowflake's task documentation covers scheduling, graphs, serverless compute, and monitoring views \cite{snowflakeTasksDocs,snowflakeAlertsDocs}. Airflow's documentation is the benchmark for external orchestration and a useful contrast when deciding what belongs in-platform \cite{airflowDocs}.

\section*{Key Takeaways}
\begin{itemize}
    \item Tasks combine schedule, compute, precondition, and body; they are created suspended and need explicit \texttt{RESUME}.
    \item Task graphs fan out and fan in through \texttt{AFTER}; \texttt{WHEN} predicates skip work---and its subtree---when no data exists.
    \item Serverless tasks fit short, frequent, spiky steps; steady heavy steps belong on controlled warehouses.
    \item Error integrations, failure suspension, and idempotent bodies turn failures into contained, pageable incidents.
    \item At 50+ transformations, layer graphs with control-table handoffs rather than one monolithic tree.
    \item \texttt{TASK\_HISTORY} is the operational source of truth; build its dashboard on day one.
\end{itemize}

\section{Exercises}
\begin{enumerate}
    \item \textbf{(Conceptual)} Explain the resume-order rule for task trees and what breaks when it is violated.
    \item \textbf{(Conceptual)} A 5-minute task regularly runs 7 minutes. What happens to the schedule, and what are two remedies?
    \item \textbf{(Conceptual)} Contrast \texttt{WHEN} gating on a parent versus on each child.
    \item \textbf{(Hands-on)} Build a fan-in graph where a report task waits on three independent merge tasks.
    \item \textbf{(Hands-on)} Convert a warehouse task to serverless and compare credits over 24 hours of 5-minute runs.
    \item \textbf{(Design)} Design the control-table schema for cross-layer handoffs, including stale-layer detection.
    \item \textbf{(Design)} Specify the backfill interface (parameters, procedures, and operator steps) for replaying a failed night.
    \item \textbf{(Operations)} Write a \texttt{TASK\_HISTORY} query listing yesterday's failures with root-cause error text and affected subtrees.
    \item \textbf{(Cost)} Estimate monthly serverless task credits for the insurer's 50-task graph given per-task durations.
    \item \textbf{(Interview)} Argue when an organization should keep Airflow instead of moving orchestration into Snowflake tasks.
\end{enumerate}

\section{Capstone Project: Real-Time CDC and Automated ELT Pipeline}\label{sec:ch8-capstone}
\index{capstone project!Milestone 2}
\begin{capstonebox}
\textbf{Mission.} Integrate Weeks 5--8 into one continuous ingestion engine: an external S3 stage with storage integration feeding auto-ingest Snowpipe into a raw VARIANT landing table; a standard stream tracking changes; and a serverless task DAG that parses JSON, executes incremental \texttt{MERGE} statements into a Gold schema, and sends error alerts on failure.

\textbf{Estimated time.} 6--9 hours across the four weeks' labs.

\textbf{Deliverable checklist.}
\begin{itemize}
    \item End-to-end ingestion runs from a clean environment following a committed README.
    \item Architecture diagram: stage $\rightarrow$ pipe $\rightarrow$ stream $\rightarrow$ task DAG $\rightarrow$ Gold, with SNS error notifications.
    \item At least three automated data-quality checks (row-count reconciliation, duplicate-key detection, JSON parse-failure quarantine) with a failing-case demonstration.
    \item Monitoring evidence: exported pipe error-integration and task error-notification configuration.
    \item Monthly cost estimate (Snowpipe credits plus serverless tasks) with two documented optimization decisions.
    \item Security review: scoped storage-integration IAM, no hardcoded secrets, least-privilege roles on stream and task objects.
\end{itemize}
\end{capstonebox}

The gating task---the capstone's hardest element---follows the pattern from Wednesday:

\begin{lstlisting}[language=SQL, caption={Serverless task gated on the stream offset}]
CREATE OR REPLACE TASK t_merge_customers  -- no WAREHOUSE = serverless
  AFTER t_gate
WHEN SYSTEM$STREAM_HAS_DATA('raw.customer_stream')
AS
  CALL sp_merge_customers();  -- idempotent MERGE procedure
\end{lstlisting}

Watch three pitfalls from the field: an unconsumed stream goes stale after 14 days (monitor \texttt{STALE} and re-snapshot); SQS events redeliver, so a non-idempotent MERGE double-applies changes; and an S3 event rule on the wrong prefix silently starves the pipe.
